%% ======================================================================
%%  Chapter III -- Noetherian rings and Dedekind rings
%% ======================================================================
\clearpage
\folio{46}%
\chapterhead{3}{III}{Noetherian rings and Dedekind rings}%
            {NOETHERIAN RINGS AND DEDEKIND RINGS}

We refer the reader who wonders why we discuss Dedekind rings to \hyperref[sec:3.4]{§4}, the
example and the discussion following Theorem \ref{thm:III.4.1}. Noetherian
rings are more general than Dedekind rings. We define Noetherian rings and
develop a few of their properties in order to place these properties in
their natural context, as well as because Noetherian rings are of
fundamental importance in other areas of algebra and in algebraic geometry.
Finally, the generalization of certain results regarding Noetherian rings to
the case of Noetherian modules is another example of ``linearization'', a
technique whose power the reader has already observed.

%% ----------------------------------------------------------------------
\sechead{1}{Noetherian rings and modules}{NOETHERIAN RINGS AND MODULES}

In \hyperref[sec:1.4]{Chapter I, §4}, we proved the following:

\begin{theorem}
\label{thm:III.1.1}
Let $A$ be a ring and $M$ an $A$-module. The following statements are
equivalent.
\begin{alist}
\item Every non-empty collection of submodules of $M$ contains a maximal
element.
\item Every increasing sequence of submodules of $M$ is stationary.
\item Every submodule of $M$ is of finite type.
\end{alist}
\end{theorem}

\begin{definition}
\label{def:III.1.1}
An $A$-module $M$ is called \emph{Noetherian} if it satisfies the equivalent
conditions of Theorem \ref{thm:III.1.1}. A ring $A$ is called
\emph{Noetherian} if, considered as an $A$-module, it is a Noetherian module.
\end{definition}

We have seen (\hyperref[sec:1.4]{Chapter I, §4}, corollary of Theorem \ref{thm:I.4.1}) that a
principal ideal ring is Noetherian.

\begin{proposition}
\label{prop:III.1.1}
Let $A$ be a ring, $E$ an $A$-module, and $E'$ a submodule of $E$. In order
that $E$ be Noetherian it is necessary and sufficient that $E'$ and $E/E'$
be Noetherian.
\end{proposition}

\begin{proof}
First, we prove necessity. Suppose $E$ is Noetherian. The lattice of
submodules of $E'$ (respectively, $E/E'$) is isomorphic to the lattice of
submodules of $E$ contained in $E'$ (respectively, containing $E'$). Thus
$E'$ and $E/E'$ are Noetherian by (a) or (b).

Conversely, suppose $E'$ and $E/E'$ are Noetherian. Let
$(F_{n})_{n\geq0}$ be an increasing sequence of submodules of $E$. As $E'$
is Noetherian, there is an integer $n_{0}$ such that
$F_{n}\cap E'=F_{n+1}\cap E'$ for all $n\geq n_{0}$. As $E/E'$ is
Noetherian, there is an integer $n_{1}$ such that
$(F_{n}+E')/E'=(F_{n+1}+E')/E'$ for all $n\geq n_{1}$. Therefore,
$F_{n}+E'=F_{n+1}+E'$ for $n\geq n_{1}$. Take $n\geq\sup(n_{0},n_{1})$. We
shall show that $F_{n}=F_{n+1}$. It suffices to show that
$F_{n+1}\subset F_{n}$. To see this take $x\in F_{n+1}$. Since
$F_{n+1}+E'=F_{n}+E'$, there exist $y\in F_{n}$ and $z',z''\in E'$ such
that $x+z'=y+z''$. Thus,
$x-y=z''-z'\in F_{n+1}\cap E'$. Note that
$F_{n+1}\cap E'=F_{n}\cap E'$. Thus, since $x-y$ and $y$ belong to $F_{n}$,
$x\in F_{n}$ too. We conclude that $F_{n+1}=F_{n}$ for all
$n\geq\sup(n_{0},n_{1})$, thus $E$ is Noetherian by (b).
\end{proof}

\begin{corollary}
\label{cor:III.1.1}
Let $A$ be a ring and let $E_{1},\ldots,E_{n}$ be Noetherian $A$-modules.
Then the $A$-module product $\prod_{i=1}^{n}E_{i}$ is Noetherian.
\end{corollary}

\begin{proof}
For $n=2$, $E_{1}$ may be identified with the submodule $E_{1}\times(0)$ of
$E_{1}\times E_{2}$, and the corresponding quotient module is isomorphic to
$E_{2}$. Our assertion follows from Proposition \ref{prop:III.1.1}. The
general case is proved by induction on $n$.
\end{proof}

\begin{corollary}
\label{cor:III.1.2}
Let $A$ be a Noetherian ring and let $E$ be an $A$-module of finite type.
Then $E$ is a Noetherian module (and, therefore, all its submodules are of
finite type).
\end{corollary}

\begin{proof}
By \hyperref[sec:1.4]{Chapter I, §4}, $E$ is isomorphic to a quotient module $A^{n}/R$ ($n$
being the cardinality of a finite set of generators of $E$). Corollary
\ref{cor:III.1.1} implies that $A^{n}$ is Noetherian and this fact, combined
with Proposition \ref{prop:III.1.1}, implies that $A^{n}/R$ is Noetherian
too.
\end{proof}

%% ----------------------------------------------------------------------
\sechead{2}{An application concerning integral elements}%
        {AN APPLICATION CONCERNING INTEGRAL ELEMENTS}

\begin{proposition}
\label{prop:III.2.1}
Let $A$ be a Noetherian integrally closed ring. Let $K$ be its field of
fractions, $L$ a finite extension of $K$, and $A'$ the integral closure of
$A$ in $L$. Suppose that $K$ is of characteristic $0$. Then $A'$ is an
$A$-module of finite type and a Noetherian ring.
\end{proposition}

\begin{proof}
We know that $A'$ is a submodule of a free $A$-module of rank $n$
(\hyperref[sec:2.7]{Chapter II, §7}, Theorem \ref{thm:II.7.1}). Thus $A'$ is an $A$-module of
finite type (\hyperref[sec:3.1]{§1}, Corollary \ref{cor:III.1.2} of Proposition
\ref{prop:III.1.1}), and, therefore, a Noetherian module (ibid.). On the
other hand, the ideals of $A'$ are special cases of $A$-submodules of $A'$.
They satisfy the maximal condition (\hyperref[sec:3.1]{§1}, Theorem \ref{thm:III.1.1}, (a)), so
$A'$ is a Noetherian ring.
\end{proof}

\begin{statement}{Example}
The ring of integers of a number field is Noetherian (take
$A=\mathbf{Z}$, $K=\mathbf{Q}$).
\end{statement}

%% ----------------------------------------------------------------------
\sechead{3}{Some preliminaries concerning ideals}%
        {SOME PRELIMINARIES CONCERNING IDEALS}

An ideal $\mathfrak{p}$ of a ring $A$ is called \emph{prime} if the quotient
ring $A/\mathfrak{p}$ is an integral domain. Equivalently, the relations
$x\in A-\mathfrak{p}$, $y\in A-\mathfrak{p}$ entail
$xy\in A-\mathfrak{p}$, i.e.\ $A-\mathfrak{p}$ is stable under
multiplication.

In order that an ideal $\mathfrak{m}$ of $A$ be \emph{maximal} (i.e.\ maximal
among the ideals of $A$ distinct from $A$), it is necessary and sufficient
that $A/\mathfrak{m}$ contain no ideals besides itself and $(0)$, i.e.\ that
$A/\mathfrak{m}$ be a field. Thus, every maximal ideal is prime. The converse
is false, as the ideal $(0)$ of $\mathbf{Z}$ is prime but not maximal.

\begin{lemma}\label{lem:III.3.1}
Let $A$ be a ring, $\mathfrak{p}$ a prime ideal of $A$, and let $A'$ be a
subring of $A$. Then $\mathfrak{p}\cap A'$ is a prime ideal of $A'$.
\end{lemma}

\begin{proof}
$\mathfrak{p}\cap A'$ is the kernel of the composition of the homomorphisms
$A'\to A\to A/\mathfrak{p}$, so there is an injective homomorphism
$A'/\mathfrak{p}\cap A'\to A/\mathfrak{p}$. Clearly, a subring of an integral
domain is an integral domain.
\end{proof}

Given two ideals $\mathfrak{a}$ and $\mathfrak{b}$ of a ring $A$, we define
the product of $\mathfrak{a}$ and $\mathfrak{b}$ not as the set of products
$ab$ where $a\in\mathfrak{a}$ and $b\in\mathfrak{b}$, but as the set of all
finite sums $\sum a_{i}b_{i}$ of such products. One sees immediately that
$\mathfrak{a}\mathfrak{b}$ is an ideal of $A$. We have:
\[
\mathfrak{a}\mathfrak{b}\subset\mathfrak{a}\cap\mathfrak{b}.
\]
The two expressions are not always equal. In a principal ideal ring the
left-hand side corresponds to the product of ideal generators and the
right-hand side to the least common multiple of generators.

Ideal multiplication is associative and commutative. $A$ acts as an identity
element in the monoid.

Given an $A$-module $E$, a submodule $F$, and an ideal $\mathfrak{a}$ of
$A$, we define in the same way the product $\mathfrak{a}F$. It is a
submodule of $E$.

\begin{lemma}\label{lem:III.3.2}
If a prime ideal $\mathfrak{p}$ of a ring $A$ contains a product
$\mathfrak{a}_{1}\mathfrak{a}_{2}\ldots\mathfrak{a}_{n}$ of ideals, then
$\mathfrak{p}$ contains at least one of the ideals $\mathfrak{a}_{i}$.
\end{lemma}

\begin{proof}
If $\mathfrak{a}_{i}\not\subset\mathfrak{p}$ for any $i$, then there exists
$a_{i}\in\mathfrak{a}_{i}-\mathfrak{p}$ for all $i$. Therefore,
$a_{1}\ldots a_{n}\notin\mathfrak{p}$, since $\mathfrak{p}$ is prime. But
$a_{1}\ldots a_{n}\in\mathfrak{a}_{1}\ldots\mathfrak{a}_{n}$, which
contradicts the hypothesis of the lemma.
\end{proof}

\begin{lemma}\label{lem:III.3.3}
In a Noetherian ring every ideal contains a product of prime ideals. In a
Noetherian integral domain $A$, every non-zero ideal contains a product of
non-zero prime ideals.
\end{lemma}

\begin{proof}
We are going to make use of a type of argument which occurs rather
frequently in the theory of Noetherian rings. Let us prove the second
assertion (the proof of the first is analogous; it suffices to delete the
word ``non-zero'' three times). We look for a contradiction. Assume that the
collection $\Phi$ of non-zero ideals of $A$ which contain no product of
non-zero prime ideals is not empty. Since $A$ is Noetherian, $\Phi$ contains
a maximal element $\mathfrak{b}$ (\hyperref[sec:3.1]{§1}, Theorem \ref{thm:III.1.1}, (a)). The
ideal $\mathfrak{b}$ cannot be prime; otherwise $\mathfrak{b}$ would not
belong to $\Phi$. Thus, there exist $x,y\in A-\mathfrak{b}$ such that
$xy\in\mathfrak{b}$. The ideals $\mathfrak{b}+Ax$ and $\mathfrak{b}+Ay$
contain $\mathfrak{b}$ as a proper subset. Therefore, since $\mathfrak{b}$ is
maximal, they do not belong to $\Phi$. It follows that they both contain
products of non-zero prime ideals:
\[
\mathfrak{b}+Ax\supset\mathfrak{p}_{1}\ldots\mathfrak{p}_{n}
\quad\text{and}\quad
\mathfrak{b}+Ay\supset\mathfrak{q}_{1}\ldots\mathfrak{q}_{r}.
\]
Since $xy\in\mathfrak{b}$,
\[
(\mathfrak{b}+Ax)(\mathfrak{b}+Ay)\subset\mathfrak{b},
\]
whence
$\mathfrak{p}_{1}\ldots\mathfrak{p}_{n}\mathfrak{q}_{1}\ldots
\mathfrak{q}_{r}\subset\mathfrak{b}$.
\end{proof}

Now let $A$ be an integral domain and let $K$ be its field of fractions. We
call any $A$-submodule $I$ of $K$ for which there exists $d\in A-(0)$ such
that $d\cdot I\subset A$ a \emph{fractional ideal} of $A$ (or of $K$ with
respect to $A$). This means that the elements of $I$ have a ``common
denominator'' $d\in A$. The ordinary ideals of $A$ are fractional ideals
(with $d=1$). We sometimes call them \emph{integral ideals} to distinguish
them from fractional ideals. Any $A$-submodule $I$ of finite type contained
in $K$ is a fractional ideal. This follows from the fact that, if
$(x_{1},\ldots,x_{n})$ is a finite set of generators for $I$, the $x_{i}$'s
have a common denominator $d$ (e.g.\ the product of the denominators $d_{i}$,
where $x_{i}=a_{i}d_{i}^{-1}$, with $a_{i},d_{i}\in A$), and $d$ is a common
denominator for $I$. Conversely, if $A$ is Noetherian, every fractional
ideal $I$ is an $A$-module of finite type, i.e.\ $I\subset d^{-1}A$ and
$d^{-1}A$, being an $A$-module isomorphic to $A$, is a Noetherian module.

We define the product $II'$ of two fractional ideals $I$ and $I'$ as the
set of finite sums $\sum x_{i}y_{i}$ where $x_{i}\in I$ and $y_{i}\in I'$.
If $I$ and $I'$ are fractional ideals, with common denominators $d$ and
$d'$, then the sets $I\cap I'$, $I+I'$, and $II'$ are all fractional ideals.
They are clearly $A$-submodules of $K$ and they have as common denominators
$d$ (or $d'$), $dd'$, and $dd'$, respectively. The non-zero fractional
ideals of $A$ constitute a commutative monoid under multiplication.

%% ----------------------------------------------------------------------
\sechead{4}{Dedekind rings}{DEDEKIND RINGS}

\begin{definition}
\label{def:III.4.1}
An integral domain $A$ is called a \emph{Dedekind ring} if it is Noetherian
and integrally closed, and if every non-zero prime ideal of $A$ is maximal.
\end{definition}

The ring $\mathbf{Z}$, and more generally any principal ideal ring, is a
Dedekind ring. The following theorem implies that the ring of integers in a
number field is a Dedekind ring.

\begin{theorem}
\label{thm:III.4.1}
Let $A$ be a Dedekind ring, $K$ its field of fractions, $L$ an extension of
finite degree of $K$, and $A'$ the integral closure of $A$ in $L$. Assume $K$
is of characteristic $0$. Then $A'$ is a Dedekind ring and an $A$-module of
finite type.
\end{theorem}

\begin{proof}
The ring $A'$ is integrally closed by construction. It is Noetherian and an
$A$-module of finite type by the proposition of \hyperref[sec:3.2]{§2}. It remains to show that
every prime ideal $\mathfrak{p}'\neq(0)$ of $A'$ is maximal. For this purpose
choose an element $x\in\mathfrak{p}'-(0)$ and consider an equation of
integral dependence of $x$ over $A$, the degree of which is a minimum:
\begin{equation}
x^{n}+a_{n-1}x^{n-1}+\cdots+a_{1}x+a_{0}=0\quad(a_{i}\in A).
\label{eq:III.4.1}
\end{equation}
Then $a_{0}\neq0$, since otherwise one could factor out an $x$ and obtain an
equation of lower degree. By \eqref{eq:III.4.1}, we have
$a_{0}\in A'x\cap A\subset\mathfrak{p}'\cap A$. Therefore,
$\mathfrak{p}'\cap A\neq(0)$. Since $\mathfrak{p}'\cap A$ is a prime ideal
of $A$ (\hyperref[sec:3.3]{§3}, \hyperref[lem:III.3.1]{Lemma 1}), we see that $\mathfrak{p}'\cap A$ is a maximal ideal of
$A$ and $A/\mathfrak{p}'\cap A$ is a field. But $A/\mathfrak{p}'\cap A$ may be
identified with a subring of $A'/\mathfrak{p}'$, and $A'/\mathfrak{p}'$ is
integral over $A/\mathfrak{p}'\cap A$ (since $A'$ is integral over $A$). Thus
$A'/\mathfrak{p}'$ is a field (\hyperref[sec:2.1]{Chapter II, §1}, Proposition \ref{prop:II.1.3}),
so $\mathfrak{p}'$ is maximal.
\end{proof}

Interest in Dedekind rings arises from the fact that the ring of integers in
a number field is a Dedekind ring, but not always a principal ideal ring.

\begin{statement}{Example}
Consider the ring of integers $A=\mathbf{Z}[\sqrt{-5}]$ in
$\mathbf{Q}[\sqrt{-5}]$ (\hyperref[sec:2.5]{Chapter II, §5}, Theorem \ref{thm:II.5.1}). Observe
that
\begin{equation}
(1+\sqrt{-5})(1-\sqrt{-5})=2\cdot3.
\label{eq:III.4.2}
\end{equation}
The norms of the four factors are, respectively, $6,6,4$, and $9$. Note that
$1+\sqrt{-5}$ can have no non-trivial divisor in $A$, since the norm of such
a divisor would have to be a non-trivial divisor of $6$. This is impossible,
because the equations $a^{2}+5b^{2}=2$ and $a^{2}+5b^{2}=3$ have no
solutions in $\mathbf{Z}$. If $A$ were principal, the prime element
$1+\sqrt{-5}$, which divides the product $2\cdot3$ by \eqref{eq:III.4.2},
would have to divide either $2$ or $3$. But then, taking norms, we see that
$6$ would divide $4$ or $9$, which is not the case.
\end{statement}

Historically, the arithmetician Kummer (1810--1893) observed that the rings
of integers in certain number fields were not principal ideal rings (in
fact, certain cyclotomic fields; Kummer observed this in connection with his
work on Fermat's equation, cf.\ I, \hyperref[sec:3.2]{§2}). In order, at least in part, to get
around this inconvenience, he and Dedekind (1831--1916) introduced the
notion of ideal. Dedekind then studied the rings which now carry his name.
The most important property of principal ideal rings is unique factorization
into products of primes. There is an elegant generalization of this property
to the case of Dedekind rings. In a Dedekind ring ideals factor uniquely
into products of prime ideals. There are many interesting consequences of
this unique factorization, which we intend now to describe precisely and
prove.

\begin{theorem}
\label{thm:III.4.2}
Let $A$ be a Dedekind ring which is not a field. Every maximal ideal of $A$
is invertible in the monoid of fractional ideals of $A$.
\end{theorem}

\begin{proof}
Let $\mathfrak{m}$ be a maximal ideal of $A$. Then
$\mathfrak{m}\neq(0)$, since $A$ is not a field. Put
\begin{equation}
\mathfrak{m}'=\{x\in K\mid x\mathfrak{m}\subset A\}.
\label{eq:III.4.3}
\end{equation}
Clearly, $\mathfrak{m}'$ is an $A$-submodule of $K$; any nonzero element of
$\mathfrak{m}$ serves as a common denominator for the elements of
$\mathfrak{m}'$. Thus $\mathfrak{m}'$ is a fractional ideal of $A$. It
suffices to show that $\mathfrak{m}\mathfrak{m}'=A$. We see that
\eqref{eq:III.4.3} implies that $\mathfrak{m}\mathfrak{m}'\subset A$; on the
other hand, $A\subset\mathfrak{m}'$ (since $\mathfrak{m}$ is an ideal), so
$\mathfrak{m}=A\mathfrak{m}\subset\mathfrak{m}'\mathfrak{m}$. As
$\mathfrak{m}$ is maximal and
$\mathfrak{m}\subset\mathfrak{m}'\mathfrak{m}\subset A$, either
$\mathfrak{m}'\mathfrak{m}=A$ or $\mathfrak{m}'\mathfrak{m}=\mathfrak{m}$. It
remains to show that $\mathfrak{m}'\mathfrak{m}=\mathfrak{m}$ is impossible.

Now, if $\mathfrak{m}'\mathfrak{m}=\mathfrak{m}$ and if
$x\in\mathfrak{m}'$, then $x\mathfrak{m}\subset\mathfrak{m}$,
$x^{2}\mathfrak{m}\subset x\mathfrak{m}\subset\mathfrak{m}$, and
$x^{n}\mathfrak{m}\subset\mathfrak{m}$ for any $n\in\mathbf{N}$ by induction.
Thus any non-zero element $d\in\mathfrak{m}$ is a common denominator for all
the powers $x^{n}$ of $x$, $n\in\mathbf{N}$. It follows that $A[x]$ is a
fractional ideal of $A$. Since $A$ is Noetherian, $A[x]$ is an $A$-module of
finite type (\hyperref[sec:3.3]{§3}, near the end), so $x$ is integral over $A$ (\hyperref[sec:2.1]{Chapter II, §1},
Theorem \ref{thm:II.1.1}). But $A$ is integrally closed; therefore,
$x\in A$; and consequently $\mathfrak{m}'\mathfrak{m}=\mathfrak{m}$ implies
$\mathfrak{m}'=A$. It remains to show that $\mathfrak{m}'=A$ is impossible.

For this purpose take a non-zero element $a\in\mathfrak{m}$. The ideal $Aa$
contains a product $\mathfrak{p}_{1}\mathfrak{p}_{2}\ldots\mathfrak{p}_{n}$
of non-zero prime ideals (\hyperref[sec:3.3]{§3}, \hyperref[lem:III.3.3]{Lemma 3}). We may take $n$ as small as
possible. We have
$\mathfrak{m}\supset Aa\supset\mathfrak{p}_{1}\ldots\mathfrak{p}_{n}$, which
means that $\mathfrak{m}\supset\mathfrak{p}_{i}$ for some $i$ (\hyperref[sec:3.3]{§3}, \hyperref[lem:III.3.2]{Lemma 2}),
say $i=1$. As $\mathfrak{p}_{i}$ is maximal by hypothesis,
$\mathfrak{m}=\mathfrak{p}_{1}$. Put
$\mathfrak{b}=\mathfrak{p}_{2}\ldots\mathfrak{p}_{n}$. Then
$Aa\supset\mathfrak{m}\mathfrak{b}$ and $Aa\not\supset\mathfrak{b}$, since
$n$ was as small as possible. There thus exists $b\in\mathfrak{b}$ such that
$b\notin Aa$. Since $\mathfrak{m}\mathfrak{b}\subset Aa$, we have
$\mathfrak{m}ba^{-1}\subset A$. According to the definition
\eqref{eq:III.4.3} of $\mathfrak{m}'$, this means that
$ba^{-1}\in\mathfrak{m}'$. But, since $b\notin Aa$, $ba^{-1}\notin A$. Thus
$\mathfrak{m}'\neq A$.
\end{proof}

\begin{theorem}
\label{thm:III.4.3}
Let $A$ be a Dedekind ring and let $P$ be the set of non-zero prime ideals
of $A$. Then
\begin{alist}
\item Every non-zero fractional ideal $\mathfrak{b}$ of $A$ may be uniquely
expressed in the form
\begin{equation}
\mathfrak{b}=\prod_{\mathfrak{p}\in P}
\mathfrak{p}^{n_{\mathfrak{p}}(\mathfrak{b})},
\label{eq:III.4.4}
\end{equation}
where, for any $\mathfrak{p}\in P$,
$n_{\mathfrak{p}}(\mathfrak{b})\in\mathbf{Z}$ and, for almost all
$\mathfrak{p}\in P$, $n_{\mathfrak{p}}(\mathfrak{b})=0$.
\item The monoid of non-zero fractional ideals of $A$ is a group.
\end{alist}
\end{theorem}

\begin{proof}
First we prove the existence of (a), i.e.\ that any fractional ideal
$\mathfrak{b}$ is a product of powers ($\geq0$ or $\leq0$) of prime ideals.
There exists $d\in A-(0)$ such that $d\mathfrak{b}\subset A$, i.e.\ such that
$d\mathfrak{b}$ is an integral ideal of $A$,
$\mathfrak{b}=(d\mathfrak{b})(Ad)^{-1}$. We may, without loss of
generality, prove (a) for integral ideals. Proceeding as in \hyperref[lem:III.3.3]{Lemma 3} of \hyperref[sec:3.3]{§3}, we
consider the collection $\Phi$ of non-zero ideals in $A$ which are not
products of prime ideals. Suppose that $\Phi$ is not empty. Let
$\mathfrak{a}$ be a maximal element of $\Phi$ ($A$ is Noetherian). Then
$\mathfrak{a}\neq A$, since $A$ is the product of the empty collection of
prime ideals. So $\mathfrak{a}$ is contained in a maximal ideal
$\mathfrak{p}$, which is thus a maximal element in the collection of
non-trivial ideals of $A$ which contain $\mathfrak{a}$. Let
$\mathfrak{p}'$ be the inverse fractional ideal of $\mathfrak{p}$. Since
$\mathfrak{a}\subset\mathfrak{p}$,
$\mathfrak{a}\mathfrak{p}'\subset\mathfrak{p}\mathfrak{p}'=A$. As
$\mathfrak{p}'\supset A$, $\mathfrak{a}\mathfrak{p}'\supset\mathfrak{a}$; in
fact $\mathfrak{a}\mathfrak{p}'\neq\mathfrak{a}$ (if
$\mathfrak{a}\mathfrak{p}'=\mathfrak{a}$ and if
$x\in\mathfrak{p}'$, then $x\mathfrak{a}\subset\mathfrak{a}$,
$x^{n}\mathfrak{a}\subset\mathfrak{a}$ for all $n$, $x$ integral over $A$,
and $x\in A$ (as in Theorem \ref{thm:III.4.2}). But this is impossible,
since $\mathfrak{p}'\neq A$ (otherwise $\mathfrak{p}'=A$ and
$\mathfrak{p}\mathfrak{p}'=\mathfrak{p}$)). According to the maximality of
$\mathfrak{a}$ in $\Phi$, we have $\mathfrak{a}\mathfrak{p}'\notin\Phi$, so
$\mathfrak{a}\mathfrak{p}'=\mathfrak{p}_{1}\ldots\mathfrak{p}_{n}$, a
product of prime ideals. Multiplying by $\mathfrak{p}$, we see that
$\mathfrak{a}=\mathfrak{p}\mathfrak{p}_{1}\ldots\mathfrak{p}_{n}$. Thus every
integral ideal of $A$ is a product of prime ideals.

Let us consider next the uniqueness of (a). Suppose that
\begin{equation}
\prod_{\mathfrak{p}\in P}\mathfrak{p}^{n(\mathfrak{p})}
=\prod_{\mathfrak{p}\in P}\mathfrak{p}^{m(\mathfrak{p})},
\quad\text{i.e.}\quad
\prod_{\mathfrak{p}\in P}\mathfrak{p}^{n(\mathfrak{p})-m(\mathfrak{p})}=A.
\label{eq:III.4.5}
\end{equation}
If $n(\mathfrak{p})-m(\mathfrak{p})\neq0$ for some prime ideals
$\mathfrak{p}\in P$, we may separate the positive and negative exponents and
write:
\[
\mathfrak{p}_{1}^{\alpha_{1}}\ldots\mathfrak{p}_{r}^{\alpha_{r}}
=\mathfrak{q}_{1}^{\beta_{1}}\ldots\mathfrak{q}_{s}^{\beta_{s}},
\]
where $\mathfrak{p}_{i},\mathfrak{q}_{j}\in P$, $\alpha_{i}>0$,
$\beta_{j}>0$, $\mathfrak{p}_{i}\neq\mathfrak{q}_{j}$ for all $i$ and $j$.
Thus $\mathfrak{p}_{1}$ contains
$\mathfrak{q}_{1}^{\beta_{1}}\ldots\mathfrak{q}_{s}^{\beta_{s}}$;
$\mathfrak{p}_{1}\supset\mathfrak{q}_{j}$ for some $j$ (\hyperref[sec:3.3]{§3}, \hyperref[lem:III.3.2]{Lemma 2}), say
$\mathfrak{p}_{1}\supset\mathfrak{q}_{1}$. But $\mathfrak{p}_{1}$ and
$\mathfrak{q}_{1}$ are both maximal, which implies
$\mathfrak{p}_{1}=\mathfrak{q}_{1}$, which is a contradiction.

Finally \eqref{eq:III.4.5} implies that
$\prod_{\mathfrak{p}\in P}\mathfrak{p}^{-n_{\mathfrak{p}}(\mathfrak{b})}$ is
the inverse of $\mathfrak{b}$ and this proves (b).
\end{proof}

\begin{statement}{Remark}
We have just seen that the monoid $I(A)$ of non-zero fractional ideals of a
Dedekind ring is a group. The principal fractional ideals (i.e.\ those of the
form $Ax$, $x\in K^{*}$) form a subgroup $F(A)$ of $I(A)$ (since
$(Ax)(Ay)^{-1}=Axy^{-1}$). The quotient group $C(A)=I(A)/F(A)$ is called the
\emph{ideal class group} of $A$. In order that $A$ be a principal ideal ring,
it is necessary and sufficient that $C(A)$ consist of a single element.
\end{statement}

Let us complete this section with some formulas, in which
$n_{\mathfrak{p}}(\mathfrak{b})$ denotes the exponent of $\mathfrak{p}$ in the
factorization of $\mathfrak{b}$ into a product of prime ideals (cf.\
\eqref{eq:III.4.4}).
\begin{equation}
n_{\mathfrak{p}}(\mathfrak{a}\mathfrak{b})
=n_{\mathfrak{p}}(\mathfrak{a})+n_{\mathfrak{p}}(\mathfrak{b})
\quad(\text{proof obvious}).
\label{eq:III.4.6}
\end{equation}
\begin{equation}
\mathfrak{b}\subset A\Leftrightarrow
n_{\mathfrak{p}}(\mathfrak{b})\geq0\quad\text{for all }
\mathfrak{p}\in P.
\label{eq:III.4.7}
\end{equation}
($\Rightarrow$ seen in the course of the proof of Theorem \ref{thm:III.4.3};
$\Leftarrow$ obvious).
\begin{equation}
\mathfrak{a}\subset\mathfrak{b}\Leftrightarrow
n_{\mathfrak{p}}(\mathfrak{a})\geq n_{\mathfrak{p}}(\mathfrak{b})
\quad\text{for all }\mathfrak{p}\in P.
\label{eq:III.4.8}
\end{equation}
($\mathfrak{a}\subset\mathfrak{b}$ means the same as
$\mathfrak{a}\mathfrak{b}^{-1}\subset A$. Now apply \eqref{eq:III.4.6} and
\eqref{eq:III.4.7}).
\begin{equation}
n_{\mathfrak{p}}(\mathfrak{a}+\mathfrak{b})
=\inf(n_{\mathfrak{p}}(\mathfrak{a}),n_{\mathfrak{p}}(\mathfrak{b})).
\label{eq:III.4.9}
\end{equation}
($\mathfrak{a}+\mathfrak{b}$ is the least upper bound of $\mathfrak{a}$ and
$\mathfrak{b}$ for ideal inclusion; to complete the proof use
\eqref{eq:III.4.8}).
\begin{equation}
n_{\mathfrak{p}}(\mathfrak{a}\cap\mathfrak{b})
=\sup(n_{\mathfrak{p}}(\mathfrak{a}),n_{\mathfrak{p}}(\mathfrak{b})).
\label{eq:III.4.10}
\end{equation}
(analogous argument involving greatest lower bounds; again use
\eqref{eq:III.4.8}).

%% ----------------------------------------------------------------------
\sechead{5}{The norm of an ideal}{THE NORM OF AN IDEAL}

In this section, $K$ denotes a number field, $n$ its degree, and $A$ the ring
of integers of $K$. We write $N(x)$ in place of $N_{K/\mathbf{Q}}(x)$.

\begin{proposition}
\label{prop:III.5.1}
If $x$ is a non-zero element of $A$, then
$|N(x)|=\operatorname{card}(A/Ax)$.
\end{proposition}

Note that, since $x\in A$, we have $N(x)\in\mathbf{Z}$ (\hyperref[sec:2.6]{Chapter II, §6},
corollary of Proposition \ref{prop:II.6.2}), so the preceding formula makes
sense.

\begin{proof}
We know that $A$ is a free $\mathbf{Z}$-module of rank $n$ (\hyperref[sec:2.8]{Chapter II, §8}),
and $Ax$ is a $\mathbf{Z}$-submodule of $A$. It is also of rank $n$, since
multiplication by $x$ maps $A$ to $Ax$ isomorphically. According to \hyperref[sec:1.5]{Chapter
I, §5}, Theorem \ref{thm:I.5.1}, there exists a base $(e_{1},\ldots,e_{n})$ of
the $\mathbf{Z}$-module $A$ together with elements $c_{i}$ of $\mathbf{N}$
such that $(c_{1}e_{1},\ldots,c_{n}e_{n})$ is a base of $Ax$.

Furthermore, the abelian group $A/Ax$ is isomorphic to the finite abelian
group $\prod_{i=1}^{n}\mathbf{Z}/c_{i}\mathbf{Z}$, whose order is
$c_{1}c_{2}\ldots c_{n}$. Write $u$ for the $\mathbf{Z}$-linear mapping of
$A$ on $Ax$ defined by $u(e_{i})=c_{i}e_{i}$ for $i=1,\ldots,n$. We have
$\det(u)=c_{1}\ldots c_{n}$. On the other hand $(xe_{1},\ldots,xe_{n})$ is
also a base for $Ax$. There is thus an automorphism $v$ of the
$\mathbf{Z}$-module $Ax$ such that $v(c_{i}e_{i})=xe_{i}$. Then $\det(v)$ is
invertible in $\mathbf{Z}$, so $\det(v)=\pm1$. But $v\cdot u$ is
multiplication by $x$, and its determinant is, by definition, $N(x)$
(\hyperref[sec:2.6]{Chapter II, §6}, \hyperref[def:II.6.1]{Definition 1}). Since $\det(v\cdot u)=\det(v)\cdot\det(u)$, we may conclude that
$N(x)=\pm c_{1}\ldots c_{n}=\pm\operatorname{card}(A/Ax)$.
\end{proof}

\begin{definition}\label{def:III.5.1}
Given a non-zero integral ideal $\mathfrak{a}$ of $A$, we call the number
$\operatorname{card}(A/\mathfrak{a})$ the \emph{norm} of $\mathfrak{a}$ and
denote it by $N(\mathfrak{a})$.
\end{definition}

Let us observe that $N(\mathfrak{a})$ is finite. In fact, if $a$ is a
non-zero element of $\mathfrak{a}$, then $Aa\subset\mathfrak{a}$, and
$A/\mathfrak{a}$ may be identified with a quotient of $A/Aa$. Thus
$\operatorname{card}(A/\mathfrak{a})\leq\operatorname{card}(A/Aa)$, which is
finite by Proposition \ref{prop:III.5.1}. On the other hand we saw that, for
a principal ideal $Ab$, $N(Ab)=|N(b)|$.

\begin{proposition}
\label{prop:III.5.2}
If $\mathfrak{a}$ and $\mathfrak{b}$ are both non-zero integral ideals of
$A$, then $N(\mathfrak{a}\mathfrak{b})
=N(\mathfrak{a})N(\mathfrak{b})$.
\end{proposition}

\begin{proof}
The ideal $\mathfrak{b}$ factors into a product of maximal ideals (\hyperref[sec:3.4]{§4},
Theorem \ref{thm:III.4.3}), and it suffices to show that
$N(\mathfrak{a}\mathfrak{m})=N(\mathfrak{a})N(\mathfrak{m})$ for
$\mathfrak{m}$ maximal. Since
$\mathfrak{a}\mathfrak{m}\subset\mathfrak{a}$, we have
$\operatorname{card}(A/\mathfrak{a}\mathfrak{m})
=\operatorname{card}(A/\mathfrak{a})
\operatorname{card}(\mathfrak{a}/\mathfrak{a}\mathfrak{m})$. It thus suffices
to show that
$\operatorname{card}(\mathfrak{a}/\mathfrak{a}\mathfrak{m})
=\operatorname{card}(A/\mathfrak{m})$. Now
$\mathfrak{a}/\mathfrak{a}\mathfrak{m}$ is an $A$-module annihilated by
$\mathfrak{m}$, which means it may be considered as a vector space over
$A/\mathfrak{m}$. Its subspaces are its $A$-submodules; and they are of the
form $\mathfrak{q}/\mathfrak{a}\mathfrak{m}$, where $\mathfrak{q}$ is an
ideal such that
$\mathfrak{a}\mathfrak{m}\subset\mathfrak{q}\subset\mathfrak{a}$. But
formula \eqref{eq:III.4.8} of \hyperref[sec:3.4]{§4} implies that there are no ideals between
$\mathfrak{a}\mathfrak{m}$ and $\mathfrak{a}$. Therefore, the vector space
$\mathfrak{a}/\mathfrak{a}\mathfrak{m}$ is of dimension one over
$A/\mathfrak{m}$. This means that
$\operatorname{card}(\mathfrak{a}/\mathfrak{a}\mathfrak{m})
=\operatorname{card}(A/\mathfrak{m})$.
\end{proof}
